\section{Security Statements and Analyses}
\label{sec:security-claims}

This section serves as the security-claim section of the submission draft. For the digital-signature functionality specified in \cref{sec:spx-core}, we claim EUF-CMA security under the assumptions and instantiations discussed below. The recommended NGCC-aligned families are recorded in \cref{tab:ngcc-parameters}, and the packaged ICCS hash/XOF interface is described in \cref{sec:hashfunctions}.

\paragraph{Theoretical security.}
Kylin follows the security-analysis framework of \SPHINCSP. As a result, the EUF-CMA security of the signature scheme is intended to reduce to the security of the tweakable hash-function families used for \hashmsg, \PRF, \PRFmsg, \hashF, \hashH, and \Th. In the usual analysis model, these primitives are treated as idealized random oracles or quantum-access random oracles, and the resulting bounds are interpreted against generic classical and quantum attacks.

In particular, the EUF-CMA security of the signature scheme relies on the following security properties of the underlying components:
\begin{itemize}
    \item The PRF security of \PRF and \PRFmsg,
    \item The interleaved target subset resilience (ITSR) of \hashmsg,
    \item The single function, multi-target undetectability (SM-UD), target-collision (SM-TCR), and decisional second-preimage resistance (SM-DSPR) of \hashF,
    \item The SM-TCR of \hashH and $\Th_\ell$.
\end{itemize}

\paragraph{Concrete security and failure probability.}

The security analysis is identical to \SPHINCSP. The classical security level, or bit security of \SPHINCSA\ against generic attacks can be computed as
$$
b=-\log\left(
    \frac{1}{2^{8n}}+
    \sum_{\gamma}
        \left(1-\left(1-\frac{1}{t}\right)^{\gamma}\right)^{k}
        \binom{q}{\gamma}
        \left(1-\frac{1}{2^{h}}\right)^{q-\gamma}
        \frac{1}{2^{h\gamma}}
    \right) \enspace.
$$
The quantum security level, or bit security of \SPHINCSA\ against generic attacks can be computed as
$$
b=-\frac{1}{2}\log\left(
    \frac{1}{2^{8n}}+
    \sum_{\gamma}
        \left(1-\left(1-\frac{1}{t}\right)^{\gamma}\right)^{k}
        \binom{q}{\gamma}
        \left(1-\frac{1}{2^{h}}\right)^{q-\gamma}
        \frac{1}{2^{h\gamma}}
    \right) \enspace.
$$
Here, we are neglecting the small constant factors inside the logarithm. 

For the recommended NGCC submission families in \cref{tab:ngcc-parameters}, we instantiate these formulas with $q=2^{64}$ in order to match the call's explicit requirement that one key pair support at least $2^{64}$ signed messages.

\paragraph{Other security properties.}
No additional formal security claims are made in this draft. In particular, the current first-stage text does not claim dedicated side-channel resistance, perfect forward security, a quantified multi-key reduction, or a separate security theorem for temporary key reuse beyond the stated single-key $q=2^{64}$ analysis.

\paragraph{Security of the ICCS hash/XOF instantiation.}

The packaged NGCC implementation uses ICCS auxiliary functions
\texttt{sm3hash}, \texttt{pseudohash}, and \texttt{pseudoXOF}. At the scheme
level, \hashmsg, \PRF, \PRFmsg, \hashF, \hashH, and \Th\ are realized through
domain-separated wrappers over \texttt{pseudoXOF} with explicit tags and
seed/address separation in the input layout.

Accordingly, the security claim in this draft is that the concrete instantiated
functions provide the pseudorandomness, target-collision resistance, and related
properties required by the \SPHINCSA\ proof framework at the selected parameter
sizes. This submission text does not introduce an additional independent
tightness claim specific to the ICCS helper internals beyond that standard
generic-analysis interpretation.

